\documentclass[11pt,a4paper]{article}
\usepackage[margin=2.4cm]{geometry}
\usepackage{amsmath}
\usepackage{booktabs}
\usepackage[hidelinks]{hyperref}
\usepackage{parskip}
\usepackage[T1]{fontenc}
\usepackage{tabularx}
\sloppy

\title{\textbf{Satisfied, Not Solved}\\[2mm]
\large What a Demographic Parity Constraint Actually Does}
\author{Dheirav}
\date{Findings summary --- \today}

\newcommand{\claim}[1]{\medskip\noindent\textbf{#1}\quad}

\begin{document}
\maketitle

\begin{abstract}
\noindent Fairness constraints of the reduction family (Agarwal et al., 2018) do what they
promise: on UCI Adult they drive the demographic parity violation from $0.186$ to $0.018$
for under two accuracy points. This work audits what they do to \emph{get} there, across
\textbf{26 populations and 61 experimental arms} spanning a household survey, its modern
replacement, and an administrative record of real mortgage decisions.

The central finding: \textbf{levelling down is not what parity constraints do --- it is what
they do below a selection rate of roughly $0.3$}. Above it, the same constraint hands
decisions out instead of taking them away. This is established by a single-factor design in
which only the label's cutoff moves, and replicates across two protected attributes, five
states, a second domain, and fourteen populations held out from a later prediction
($r = +0.901$).
\end{abstract}

\section{The problem in one paragraph}

A constraint that equalises two groups' selection rates can be satisfied two ways: by
extending favourable decisions to the disadvantaged group, or by withdrawing them from the
advantaged one. Every real mitigation does a mixture of both. \textbf{The certifying metric
is identical either way.} It reports that the violation fell; it says nothing about which
mixture produced that, or whether the total pool of favourable decisions grew or shrank.

\section{What is already known, and what is not}

\claim{Known.} Mittelstadt, Wachter and Russell (2023) established that fairness
interventions frequently equalise by degrading the better-off group, and argue it is a
default rather than an accident. Their \S6 also proposes the remedy --- ``minimum rate
constraints'', requiring that every group has at least a minimal selection rate ---
demonstrated on Adult, and reporting that ``levelling down does not occur''.

\claim{Therefore not claimed here.} Neither the observation nor the remedy is novel, and
this document says so before a reviewer does. An earlier version of this project claimed
the remedy as new; that was wrong about the literature and has been corrected in place.

\claim{Not known.} \emph{When} it happens. Mittelstadt et al.\ treat levelling down as a
default; their paper contains no discussion of base rates. Their only remark touching
prevalence is that Adult is ``more than 75\% negatively labelled'', offered as a note on
accuracy cost and never connected to direction. They were operating at a selection rate of
roughly $0.24$ --- inside the levelling-down regime identified here --- and did not ask
whether it would reverse elsewhere.

\claim{The nearest miss.} Goethals, Delaney, Mittelstadt and Russell (2024),
\emph{Resource-constrained Fairness}, study the \emph{cost} of fairness as a function of the
available budget and note the selection rate is ``a factor overlooked in previous
evaluations''. Two of the four authors are also authors of the levelling-down paper. But in
their setting positive decisions are a \emph{fixed} resource, so the total cannot move and
the directional effect reported here is unobservable by construction. The group best placed
to find this result studied the same axis and did not connect it to direction.

\section{The central result}

\subsection{The design}

ACS Income's label is ``earns more than \$50{,}000'' --- a cutoff chosen by the benchmark,
not a fact about the world. Moving it on a \emph{fixed} population varies the overall
selection rate while holding the population, the survey instrument, the feature set, the
group ratio and the proxy structure exactly fixed. Test code asserts that rows, features and
protected groups are identical across arms and only the label changes.

\textbf{The cutoff is the instrument, not the cause.} The moderator is the selection rate
itself. This is confirmed in a domain where no cutoff exists: HMDA mortgage decisions, where
the label is a real lender's approve/deny, sits naturally at a selection rate of $0.808$ and
lands exactly where the sweep predicts.

\subsection{Alabama, six cutoffs, five seeds}

\begin{center}
\begin{tabular}{lrrr}
\toprule
Cutoff & Selection rate & Change in favourable decisions & Destroyed per created \\
\midrule
\$100{,}000 & 0.030 & $-29.71\%$ & 22.03 \\
\$70{,}000  & 0.099 & $-22.05\%$ & 2.18 \\
\$50{,}000  & 0.252 & $-2.34\%$  & 1.14 \\
\$30{,}000  & 0.598 & $+0.98\%$  & 0.83 \\
\$20{,}000  & 0.760 & $+0.80\%$  & 0.75 \\
\$10{,}000  & 0.890 & $+0.08\%$  & 0.89 \\
\bottomrule
\end{tabular}
\end{center}

Same people. Same features. Same groups. Only the finish line moved --- and the direction
flipped, with the crossover between $0.25$ and $0.60$.

\subsection{Replication}

\begin{center}
\begin{tabular}{lll}
\toprule
Evidence & Scope & Result \\
\midrule
Original sweep & Alabama, Oregon (sex) & $r = +0.801$, $+0.964$ \\
Second protected attribute & 4 states (race) & $r = +0.874$ to $+0.991$ \\
Held-out populations & 14, run after prediction & $r = +0.901$ \\
Second domain & HMDA, 2 states, 2 arms & correct side of crossover \\
Datasets never fitted to & Adult ($0.205$), HMDA ($0.808$) & both placed correctly \\
\bottomrule
\end{tabular}
\end{center}

\subsection{The obvious objection, pre-empted}

That unequal base rates \emph{between groups} force trade-offs is textbook (Kleinberg et
al., 2016; Chouldechova, 2017). \textbf{This claim is about a different quantity}: the
overall \emph{level} of the selection rate, irrespective of any gap.

The two are confounded by construction --- the between-group gap must vanish as the level
approaches $0$ or $1$. So the gap is partialled out, and the relationship
\emph{strengthens}: from $r = +0.801$ to $+0.980$ in Alabama, and $+0.964$ to $+0.994$ in
Oregon. That control is the identification strategy, not a robustness check.

\section{An identity that says precisely what levelling down is}

Let the privileged group be share $p$ with unconstrained selection rate $r_A$, and the
unprivileged group $1-p$ at $r_B$. The unconstrained total is the size-weighted average
$\bar{r} = p\,r_A + (1-p)\,r_B$. Demographic parity puts \emph{both} groups on one common
rate $s$, so the constrained total \emph{is} $s$, and
\[
\text{change in favourable decisions} \;=\; s - \bar{r}.
\]
Writing $\lambda = (s - r_B)/(r_A - r_B)$ for where the common rate falls between the two
original rates,
\[
\boxed{\;\text{the pool is preserved exactly when } \lambda = p.\;}
\]
Levelling down is $\lambda < p$ and nothing else: the optimiser choosing a compromise
\emph{below} the size-weighted one. This requires no assumptions, holds to $0.076$
percentage points across fourteen held-out populations, and does not appear to be stated
anywhere in the literature reviewed.

\section{What could not be established}

\claim{The mechanism.} A derivation was attempted: under group-wise thresholding,
calibration and a location shift, levelling down is a \emph{curvature} effect and the
crossover should sit at the mode of the score density. It was pre-registered with numerical
thresholds and tested on fourteen populations run afterwards. It cleared its bars
(predicting the sign in 12 of 14) and was \textbf{beaten by the constant rule ``levelling up
iff the selection rate exceeds $0.5$'' (13 of 14)}. The measured mode moves inversely with
the selection rate ($r = -0.802$), so it carries almost no independent information.

The methodological lesson is recorded with it: a \emph{threshold} was fixed in advance, but
a \emph{baseline to beat} was not, so clearing the threshold did not mean what it appeared
to mean.

\claim{The magnitude.} The selection rate sets the direction, not the size. Alabama loses
$2.34\%$ at a rate of $0.252$ where Adult loses $20.5\%$ at $0.205$. A pre-registered
candidate --- the group ratio --- was refuted, with the opposite sign to the prediction.

\section{Supporting findings}

\claim{Subgroups.} A constraint on one attribute leaves Sex\,$\times$\,Race subgroups worse
off than any group it was told about, in every sufficiently diverse population, and more
severely than Adult suggests ($9.0\times$ against $13.2\times$). In five of ten populations
the worst-off subgroup after a sex constraint is a minority man.

\claim{Attribution audits do not reveal the mechanism.} A parity constraint moves SHAP
attribution \emph{onto} a sex proxy by $+151\%$. Three explanations were proposed; two were
refuted by intervention and the third bounded by re-aggregation --- scoring the two most
redundant features as one coalition takes the effect from $+155\%$ to $+11.6\%$. The
transferable lesson: \textbf{an attribution audit can move by an order of magnitude without
the model's behaviour changing correspondingly.}

\claim{Arbitrariness.} Below roughly $2{,}500$ test subjects, the method's own randomness
exceeds the entire effect of the constraint.

\claim{Against existing remedies.} Group-wise thresholding --- the \emph{optimal}
DP-constrained classifier --- destroys $18.8\%$ of favourable decisions on Adult against the
reduction's $20.5\%$, so levelling down is not a solver artifact. Minimax group fairness,
proposed in the literature as \emph{the} remedy for levelling down, destroys $10.0\%$ of
favourable decisions on HMDA at an exchange rate of $46.7$ --- on the same data where the
parity constraint \emph{creates} $4.3\%$. Its objective minimises the worst group's
\emph{error}, and worst-off-by-error is not worst-off-by-outcome: on Adult the higher-error
group is Male.

\section{Why it matters}

Lending, hiring and admissions --- the domains where these methods are actually proposed ---
are nearly all low-selection-rate settings. \textbf{The default regime for deployed fairness
constraints is therefore the one in which they take favourable decisions away, and the
fairness score never says so.}

\section{Method, and what it caught}

Five seeds throughout; splits stratified on the (attribute, label) interaction; the
protected attribute removed from the feature matrix; metrics implemented from their
definitions and cross-checked against \texttt{fairlearn} on every run; undefined rates
returned as undefined rather than zero.

Where a result could have gone either way, predictions and their numerical thresholds were
written and committed \emph{before} the experiment ran, verifiable in the commit history.
That discipline is why the following can be reported: \textbf{three pre-registered
predictions of my own failed}; two proposed mechanisms were refuted by my own interventions;
and three claims were corrected or retracted in place, including the novelty claim for the
remedy, which the literature had anticipated.

\section{The open question}

Is the scope condition sufficient on its own for a venue such as FAccT or AIES, or should
the mechanism be established first? The derivation attempted here failed, and is reported as
a failure. The finding is \emph{when} the direction flips, not \emph{why}.

\vfill
\noindent\rule{\textwidth}{0.4pt}\\
{\small All figures are reproduced from stored results by
\texttt{tests/test\_documented\_claims.py}, which re-derives every headline number and fails
if a document and its data disagree. Supporting papers are in \texttt{papers/}; the full
research record is in \texttt{research/docs/}.}

\end{document}
